\documentclass[11pt,a4paper]{article}
\usepackage[margin=24mm]{geometry}
\usepackage{amsmath,amssymb,array}
\input{glyphtounicode}
\pdfgentounicode=1
\usepackage[hidelinks]{hyperref}
\hypersetup{hypertexnames=false}
\setlength{\parindent}{0pt}
\setlength{\parskip}{5pt}
\newcommand{\pos}[1]{[#1]_+}
\begin{document}
\begin{center}
{\Large Revised A1-C: service work and causal traffic state}\\[4pt]
15 October 2026 --- conditional module, diagnostic and new choice
\end{center}

\textbf{Status.} Revised A1-C's qualification direction is user-approved;
M2-A1 is closed at that scope. The equations below are proposed/diagnostic,
not an approved complete M2 traffic law. Stop at A1-C-S1 before selecting
the state/speed dynamics. No simulator or traffic experiment is involved.

\section*{1. One work account and a concrete allocation module}

Retain the finite terminal histories, common advice, triangular backbone
and physical/eligible stocks:
\begin{equation}
 P_j(t)=F_j(t-b/v),\quad q_j=P_j-Y_j,\quad N_j=F_j-Y_j,\quad
 C(v)=\frac{wk_jv}{w+v},\qquad j=A,E,B. \tag{1}
\end{equation}
As before, eligible demand is $d_j=C(v)$ for $q_j>0$, and
$d_j=\dot P_j$ at an empty eligible queue. All lane offers are at most
$C(v)$. A causal solution must maintain $Y_j\leq P_j$.

Let $X$ denote the justified current aggregate traffic state and $Z$ the
retained service/recovery state. For fixed inputs, consider
\begin{equation}
 0\leq B\leq R\leq C(v_-),\qquad c_A(X,Z),c_E(X,Z)\geq1,
 \qquad c_Ag_A+c_Eg_E\leq B. \tag{2}
\end{equation}
$B$ is available normalized service work per second; $g_j$ is admitted
vehicle-equivalent mass per second; $c_j$ is work per unit mass relative
to ordinary passage. Only when $c_j=1$ do work and mass rates coincide.
There is one budget, not a separate recovery constraint per past entry.
The complete model must predict $B$ from its account, not receive a future
service trace from outside. Here it is an input only to the local module.
The maps $B,c_j$ and the physical division of current versus retained
burden remain to be qualified; they cannot both charge the same effect.

\textbf{Proposed fair-work convention.} For these fixed inputs choose the
smallest $z\geq0$ with
\begin{equation}
 H(z):=\min\{c_Ad_A,z\}+\min\{c_Ed_E,z\}
     =\min\{B,c_Ad_A+c_Ed_E\},\quad
 g_j=\min\{d_j,z/c_j\}. \tag{3}
\end{equation}
If both offers are zero, take $z=0$. $H$ is continuous and strictly
increasing until all offered work is served; hence (3) determines unique
flows. If work demand does not exceed $B$, both offers pass. Otherwise
the full work budget is used, with unused shares reallocated. In all cases
\begin{equation}
 0\leq g_j\leq d_j,\qquad g_A+g_E\leq B\leq R. \tag{4}
\end{equation}
Equation (3) is a declared equal-work fairness assumption, not a theorem
of traffic behavior. Conservative allocation templates motivate its shape;
the source literature does not establish our state-to-work maps.

\newpage
\section*{2. Shares, effective throughput and residual service}

When both streams can use their work shares,
\begin{equation}
 z=B/2,\qquad g_A=\frac{B}{2c_A},\quad g_E=\frac{B}{2c_E},\quad
 s_E=\frac{g_E}{g_A+g_E}=\frac{c_A}{c_A+c_E}. \tag{5}
\end{equation}
Thus $c_E>c_A$ gives E a smaller mass share and also changes total flow.
For $c_A=c_E=c$, the shares are equal but total service is $B/c$.
For $c_E=c_A(1+\varepsilon)$, E's share is $1/(2+\varepsilon)$.
These are conditional implications, not an approved equality or small
$\varepsilon$ assumption for real lane-drop traffic.

For an actual or specified offered mix $\theta_j$ summing to one,
\begin{equation}
 \mu_{\rm mix}=\frac{B}{\theta_Ac_A+\theta_Ec_E}. \tag{6}
\end{equation}
At binding work demand, actual total outflow equals the reading for its
admitted mix; at underload it may be smaller. Do not infer capacity loss
from low outflow. A composition-independent capacity is not asserted by
(6), and no mix is inferred from absent traffic. This is one shared
account with a composition-dependent mass conversion.

If A's offer is served fully and E remains work-limited,
\begin{equation}
 g_A=d_A,\qquad g_E=\frac{B-c_Ad_A}{c_E},\qquad
 c_Ad_A\leq B/2,\quad c_Ed_E\geq B-c_Ad_A. \tag{7}
\end{equation}
A larger E cost can reduce use of nominal spare mass service, unlike a
priority weight that only redistributes a fixed mass budget.

\begin{center}
\begin{tabular}{c c c c c}
$B$ & $c_A$ & $c_E$ & $(g_A,g_E)$ when both ready & total\\\hline
6 & 1 & 1 & $(3,3)$ & 6\\
6 & 1 & $3/2$ & $(3,2)$ & 5\\
6 & 1 & 2 & $(3,3/2)$ & $9/2$\\
6 & 2 & 2 & $(3/2,3/2)$ & 3
\end{tabular}
\end{center}

These normalized exact hand allocations do not assign physical traffic
parameters. For $B=6,d_A=2,c_A=1,c_E=2$ and ample E offer, (7) gives
$(g_A,g_E)=(2,2)$. At underload, e.g. $d_A=d_E=1$ with these costs,
both pass and unused work is positive.

\textbf{Not yet a full model.} Fixed-input uniqueness does not establish
uniqueness after costs depend on endogenous speed, queues or recovery.
The next page tests one tempting instantaneous closure. No numerical
smoothing, clipping epsilon or timestep-selected branch is introduced.

\newpage
\section*{3. An exact queue-birth ambiguity in a simple completion}

Use normalized $v=w=1,k_j=2$, so $C(v)=1$ and the congested FD speed is
$u_{\rm c}(f)=f/(2-f)$. Over a local time interval take
\begin{equation}
 B=R=4/5,\quad \dot P_A=a=1/5,\quad \dot P_E=\lambda=1/2,
 \quad q_A=0,\quad q_E(0)=0,\quad c_A=1. \tag{8}
\end{equation}
The combined ordinary offer is $7/10<4/5$. Illustrate a monotone speed-cost
law and a queue-selected speed by
\begin{equation}
 c_E=2-u_E,\qquad
 u_E=\begin{cases}1,&q_E=0,\\ g_E/(2-g_E),&q_E>0.\end{cases} \tag{9}
\end{equation}
This cost law is a diagnostic assumption, not an empirical formula or
the selected A1-C closure. The recovery budget is held fixed to isolate
the instantaneous entry-state test. This local constant-offer interval
is not asserted to be an exact H1-generated finite-cohort MP6 witness.

At $q_E=0$, $c_E=1$ and work demand $a+\lambda=7/10<B$, so
$g_A=a,g_E=\lambda$ and $\dot q_E=0$. At $q_E>0$, E can send at 1
and A uses its full offer. Equation (7) requires
\begin{equation}
 g_E\left(2-\frac{g_E}{2-g_E}\right)=\frac35
 \quad\Longleftrightarrow\quad15g_E^2-23g_E+6=0. \tag{10}
\end{equation}
The roots are $1/3$ and $6/5$; only $1/3$ is feasible. Then
$u_E=1/5,c_E=9/5,g_A=1/5$, the work budget is exactly $4/5$, and
\begin{equation}
 \dot q_E=\lambda-g_E=1/6\quad(q_E>0),\qquad
 \dot q_E=0\quad(q_E=0). \tag{11}
\end{equation}
For every chosen $t_b\geq0$, as well as $t_b=\infty$, the function
\begin{equation}
 q_{t_b}(t)=\frac{\pos{t-t_b}}6,\quad
 Y_E(t)=t/2-q_{t_b}(t),\quad Y_A(t)=t/5 \tag{12}
\end{equation}
satisfies the absolutely-continuous queue/departure rules almost everywhere
on the local offer interval. Both departures are nonnegative, increasing,
eligible and budget-feasible. Yet queue birth is arbitrary.

The differential equation holds outside isolated transition times and
does not choose the next regime there. An explicit rule for continuation
could exclude (12), but it would be an additional scientific rule. A
discrete integrator selecting a branch does not qualify that rule. This
example concerns the stated closure, not all A1-C laws or a complete
P-H history model.

\newpage
\section*{4. Conditional physical and clearance interfaces}

If a complete causal law supplies admissible departures, retain
\begin{equation}
 H=k_jb>M,\qquad F_j\leq M,\qquad 0\leq\dot F_j\leq C(v). \tag{13}
\end{equation}
For the initially empty terminal, vacancy before $b/w$ is positive since
$F_j(t)\leq C(v)t<k_jwt$. Afterwards
\begin{equation}
 V_j(t)=Y_j(t-b/w)+H-F_j(t)\geq H-M>0. \tag{14}
\end{equation}
Thus the sufficient domain protects the H1 handoff under either future
state closure. This is not a general upstream spillback solver.

Physical stock is $N_j=F_j-Y_j$, never service memory or weighted mass.
For uniform initial footprint $[0,\ell]$ and road length $L$, the complete
local cohort objective remains
\begin{equation}
 J(v)=\frac{M(L-\ell/2)}v+\int_0^\infty(q_A+q_E+q_B)\,dt. \tag{15}
\end{equation}
Actual departure mass is $g_jdt$, not $c_jg_jdt$. A completed entry may
retain service influence while contributing no further local vehicle-time.
B retains its independent admitted domain and $Y_B=P_B$.

\textbf{A sufficient clearance condition only.} All traffic is eligible
by $t_*=L/v$, and B is clear then. Suppose the future qualified account
has uniform bounds
\begin{equation}
 1\leq c_j(t)\leq c_{\max}<\infty,\quad 0\leq B(t)\leq R,
 \quad\int_{t_*}^\infty(R-B(t))dt\leq\Lambda<\infty. \tag{16}
\end{equation}
If $Q=q_A+q_E>0$, one lane offers $C(v)\geq R$, hence its work
offer reaches at least $B$. The account binds and
$B=c_Ag_A+c_Eg_E\leq c_{\max}(g_A+g_E)$. Before clearance,
\begin{equation}
 Q(T)\leq Q(t_*)-\frac{R(T-t_*)-\Lambda}{c_{\max}}. \tag{17}
\end{equation}
Therefore clearance occurs by
\begin{equation}
 t_*+\frac{c_{\max}Q(t_*)+\Lambda}{R}. \tag{18}
\end{equation}
Temporary zero budget is allowed; no artificial positive floor is assumed.
The complete M2 law must establish (16) or a different valid clearance
argument. The work module alone and M1's finite-kernel proof do not do so.

\textbf{Information and gradients.} Known initialization plus a uniquely
autonomous state law could preserve ending-statistic reconstruction.
That is conditional: no full-law uniqueness, consequential reachable MP6
regime, observation collision, action split or useful gradient is proved.
Changes in cost, state and event times all enter the full-objective path;
none may be frozen merely to simplify a later derivative check.

\newpage
\section*{5. What must be selected before exact-law qualification}

\textbf{S-E: equilibrium boundary state plus explicit causal transitions.}
The free branch persists while its physical admission/service conditions
permit it. A specified restriction must trigger queue formation; a specified
clearance condition returns the stream to free passage. This rejects the
arbitrary births in (12) if the free conditions remain satisfied. It does
not supply a complete entropy/transition rule or prove the whole history
node unique. A simple speed-only cost can remain benign under ordinary
underload; conditional compulsory entry/opportunity must still be modeled.

\textbf{S-V: finite-rate local aggregate speed state.}
Predict approach/entry speeds from explicit local initial-value dynamics
and admission/empty-state rules. Queue presence is a constraint/input,
not an instantaneous complete speed assignment. This changes local
motion/state meaning and needs parameter and averaging semantics. It is
not automatically microscopic car following or a spatial CTM model, and
an added speed ODE alone does not guarantee a valid node.

Both choices may use a unified service map that has intermediate states;
neither requires linear or globally smooth interpolation. They differ
in how the traffic state changes through time, not in whether forced
interaction and ordinary congestion may coexist. Neither is selected.
The provisional first qualification is S-E because it retains the smaller
first-order amendment; its failure would warrant an explicit return to S-V.

\textbf{Sources and adaptation.} The boundary formulation in
\href{https://arxiv.org/pdf/1310.2660}{Jin et al.} (sections 2 and 6) has
its own selection rule and proof. The inspected relation in
\href{https://repository.tudelft.nl/file/File_c424c1b3-7374-436b-a6a7-abdf69063aa6}{Leclercq et al.}
(section 3.2) assumes congested entry. Neither paper supplies this work
node. The fairness convention, state-to-cost proposal and queue-birth
diagnostic are review constructions, not proofs of the finite task.

\textbf{Gate and original-stage status.} A1-C-S1 is Type 2 G3; select the
local state/dynamics interpretation before qualifying its exact equations,
recovery/overlap and joint MP6 domain. M2-A1's revised A1-C direction is
closed. H1/cutoff, known initial information, finite storage, primary
ordinary-underload domain, local choke completion and S-P--S-L remain.
The approved three-part review and Step 3 are unfinished. No simulator,
installation, calibration, traffic experiment or numerical gradient has
been undertaken. The strict, fixed-share, aggregate/history/scalar and
search rivals remain; a common fastest action is an admissible outcome.
\end{document}
